\scanpage{40}
\inserted{that $\operatorname{card}I<\operatorname{card}J-1$ i.e. $\operatorname{card}(J\setminus I)=\operatorname{card}(Y\setminus U)>1$,}
\deleted{We are therefore in one of the \unsure{following|following} cases}

\deleted{a) $f^{-1}(I)=I$ i.e. $I$ stable under $f$.}

\deleted{b)} Thus we see that we must have $f^{-1}(I)=I$ i.e. $I$ stable under $f$, which shows that $\tau_u$ in fact induces an \uline{automorphism} of $Y$.
\margin{\inserted{Replacing} \deleted{\unsure{We therefore have}}}

\deleted{that $Y$ is an elliptic curve}

in these arguments $Y=\widehat Y\setminus I$ by $Y'=Y\cup i=\widehat Y\setminus I'$ $(I'=I-\{i\})$, we see likewise that $f$ induces an automorphism of $Y'$ (this argument works as soon as $Y'$ injects into its generalized Jacobian, which \deleted{except in the case where $Y$ \unsure{is} of type $(0,\nu)$ \unsure{with}} \inserted{is the case thanks to the anabelian hypothesis on $Y$)}

\deleted{\unsure{and|that is}} thus $f$ induces the identity on $I$. If the genus is $0$, we conclude from this (since $\nu\geq3$) that $f$ is the identity; \deleted{There remains the} In the case of genus $1$, we now conclude from this that the image of $u$ in $J_{\widehat Y}^0$ is equal to $1$, and we finish as before.

Phew!
\margin{39}
